\section{Conclusion}\label{sec:conclusion}

We have identified and solved a fundamental limitation of phase-field Allen-Cahn LBM for simulating droplet dynamics on curved surfaces. The Allen-Cahn sharpening term resists geometric wetting boundary condition corrections, suppressing contact angle accuracy on curved hydrophobic surfaces. We proposed a geometric amplification factor $\alpha_{geo}$ that overcomes this resistance, derived it from the balance between AC sharpening and geometric wetting, and systematically calibrated it across three parameter dimensions.

The key finding is that geometric amplification is effective only for strongly hydrophobic surfaces ($\theta_{eq} > 120°$), small curvature ratios ($R^* < 3$), and sufficient interface resolution ($D_0/\xi > 3.0$). The optimal $\alpha_{geo} = 1.5$ is recommended as a universal default for superhydrophobic curved surfaces, reducing the contact angle error from $\sim 13°$ to $\sim 3°$. Using this correction, we constructed the first We $\times$ $R^*$ phase diagram of asymmetric droplet spreading via Allen-Cahn LBM at 828:1 density ratio, revealing three spreading regimes with the asymmetry ratio $k$ peaking at $\mathrm{We} \approx 15$ ($k = 3.30$ at $R^*=1.0$).

The method has been validated across six curvature ratios ($R^* = 0.5$--$3.0$), four contact angles ($\theta_{eq} = 90°$--$162°$), and five grid resolutions ($N = 80$--$180$), with quantitative agreement within $2.1\%$ of the experimental benchmark of Liu et al.~\cite{liu2015symmetry} at $R^*=1.0$. The method correctly preserves symmetric spreading on concave surfaces ($k = 1.0$), where no asymmetry is physically expected. Grid convergence is achieved at $N \geq 100$, with the asymptotic value $k_\infty \approx 2.66 \pm 0.04$.

Limitations include: (1) quantitative validation is limited to the cylindrical ridge geometry of Liu et al.; (2) the contact time measurements are approximate and have not been compared with experimental data; and (3) the amplification factor is held constant rather than adapting to local curvature. Promising future directions include extending the calibration to three-dimensional hemispherical and concave geometries, developing an adaptive $\alpha_{geo}$ that varies with local curvature and interface thickness, coupling with thermal models for icing applications on curved surfaces, and constructing machine learning surrogates trained on the phase diagram data for real-time prediction of droplet dynamics.
